% ===========================================================================
%  PREAMBULE
%  Paquets volontairement limites a des paquets standards disponibles
%  sur toute distribution TeX Live / MiKTeX recente et sur Overleaf.
% ===========================================================================

% ---- Encodage et langue ---------------------------------------------------
\usepackage[utf8]{inputenc}
\usepackage[T1]{fontenc}
\usepackage[french]{babel}
\usepackage{lmodern}

% ---- Mise en page ---------------------------------------------------------
\usepackage[a4paper,margin=2.5cm,headheight=14pt]{geometry}
\usepackage{fancyhdr}
\usepackage{setspace}
\onehalfspacing

% ---- Tableaux et figures --------------------------------------------------
\usepackage{graphicx}
\usepackage{booktabs}
\usepackage{longtable}
\usepackage{array}
\usepackage{tabularx}
\usepackage{multirow}
\usepackage{float}
\usepackage{caption}

% ---- Couleurs et encadres -------------------------------------------------
\usepackage{xcolor}
\usepackage{framed}

% ---- Listes ---------------------------------------------------------------
\usepackage{enumitem}

% ---- Code source ----------------------------------------------------------
\usepackage{listings}

% ---- Symboles -------------------------------------------------------------
\usepackage{amssymb}
\usepackage{pifont}

% ---- Liens hypertexte (a charger en dernier) ------------------------------
\usepackage[hidelinks,breaklinks=true]{hyperref}

% ===========================================================================
%  COULEURS DU DOCUMENT
% ===========================================================================
\definecolor{bleufonce}{RGB}{11,61,92}
\definecolor{bleuclair}{RGB}{244,247,249}
\definecolor{grisbord}{RGB}{180,180,180}
\definecolor{orangeattention}{RGB}{217,145,0}
\definecolor{jauneclair}{RGB}{255,248,230}
\definecolor{vertfonce}{RGB}{45,107,45}
\definecolor{vertclair}{RGB}{234,244,234}
\definecolor{rougeplaceholder}{RGB}{176,0,32}
\definecolor{griscode}{RGB}{245,245,245}
\definecolor{commentairevert}{RGB}{0,110,0}
\definecolor{chainemauve}{RGB}{140,0,140}

% ===========================================================================
%  EN-TETES ET PIEDS DE PAGE
% ===========================================================================
\pagestyle{fancy}
\fancyhf{}
\fancyhead[L]{\small\nouppercase{\leftmark}}
\fancyhead[R]{\small\thepage}
\renewcommand{\headrulewidth}{0.4pt}

\fancypagestyle{plain}{%
  \fancyhf{}
  \fancyhead[R]{\small\thepage}
  \renewcommand{\headrulewidth}{0pt}
}

% ===========================================================================
%  STYLE DU CODE SOURCE
% ===========================================================================
\lstdefinestyle{stylegeneral}{
  backgroundcolor=\color{griscode},
  basicstyle=\ttfamily\scriptsize,
  keywordstyle=\color{bleufonce}\bfseries,
  commentstyle=\color{commentairevert}\itshape,
  stringstyle=\color{chainemauve},
  numberstyle=\tiny\color{gray},
  numbers=left,
  numbersep=8pt,
  frame=single,
  rulecolor=\color{grisbord},
  breaklines=true,
  breakatwhitespace=false,
  showstringspaces=false,
  tabsize=2,
  captionpos=b,
  xleftmargin=14pt,
  framexleftmargin=14pt,
  literate=%
    {é}{{\'e}}1 {è}{{\`e}}1 {ê}{{\^e}}1 {ë}{{\"e}}1
    {à}{{\`a}}1 {â}{{\^a}}1 {ô}{{\^o}}1 {î}{{\^i}}1
    {ï}{{\"i}}1 {ù}{{\`u}}1 {û}{{\^u}}1 {ç}{{\c c}}1
    {É}{{\'E}}1 {È}{{\`E}}1 {À}{{\`A}}1
}
\lstset{style=stylegeneral}

\lstdefinelanguage{yamlperso}{
  keywords={name,on,jobs,steps,uses,with,run,runs-on,permissions,
            concurrency,if,needs,env,repos,hooks,rev,id,args},
  sensitive=true,
  comment=[l]{\#},
  morestring=[b]",
  morestring=[b]'
}

% ===========================================================================
%  COMMANDES PERSONNALISEES
% ===========================================================================

% --- Reference a une pratique DevSecOps : \pratique{P1}{SAST}
\newcommand{\pratique}[2]{\textbf{#1}~---~#2}

% --- Identifiant de pratique seul : \pid{P1}
\newcommand{\pid}[1]{\textbf{#1}}

% --- Exigence ASVS : \asvs{8.2.2}
\newcommand{\asvs}[1]{\mbox{\textsc{asvs}~v5.0.0--#1}}

% --- Exigence ASVS courte (dans un tableau) : \asvsc{8.2.2}
\newcommand{\asvsc}[1]{\mbox{\textsc{asvs}~#1}}

% --- Risque OWASP LLM : \llm{01}
\newcommand{\llm}[1]{\mbox{\textsc{llm}#1}}

% --- Marqueurs d'etat
\newcommand{\fait}{\textcolor{vertfonce}{\ding{51}}}
\newcommand{\nonfait}{\textcolor{grisbord}{\ding{55}}}
\newcommand{\encours}{\textcolor{orangeattention}{$\circlearrowright$}}
\newcommand{\bloque}{\textcolor{rougeplaceholder}{\ding{47}}}

% ===========================================================================
%  EMPLACEMENT DE CAPTURE D'ECRAN MANQUANTE
%  Affiche un encadre bien visible dans le PDF pour ne rien oublier.
%  Usage : \captureaplacer{figures/nom.png}{description de la capture}
% ===========================================================================
\newcommand{\captureaplacer}[2]{%
  \begin{figure}[H]
  \centering
  \setlength{\fboxrule}{1.2pt}%
  \setlength{\fboxsep}{10pt}%
  \textcolor{rougeplaceholder}{%
    \fbox{%
      \begin{minipage}{0.86\textwidth}
        \centering
        \color{rougeplaceholder}%
        \textbf{\large CAPTURE D'ÉCRAN À AJOUTER}\\[6pt]
        \color{black}%
        \textbf{Fichier attendu :} \texttt{#1}\\[4pt]
        \textbf{Contenu :} #2\\[8pt]
        {\small\itshape Déposer l'image dans \texttt{figures/} sous ce nom
        exact, puis remplacer \texttt{\textbackslash captureaplacer}
        par \texttt{\textbackslash captureprete} dans le fichier source.}
        \vspace{4pt}
      \end{minipage}%
    }%
  }
  \end{figure}
}

% ===========================================================================
%  CAPTURE FOURNIE
%  Usage : \captureprete{figures/nom.png}{Légende de la figure}
%
%  ROBUSTE PAR CONCEPTION : si le fichier image est absent, la commande
%  affiche l'encadre de rappel au lieu de provoquer une erreur de
%  compilation. Un fichier oublie ne peut donc pas casser le document.
% ===========================================================================
\newcommand{\captureprete}[2]{%
  \IfFileExists{#1}{%
    \begin{figure}[H]
      \centering
      \includegraphics[width=0.92\textwidth]{#1}
      \caption{#2}
    \end{figure}
  }{%
    \captureaplacer{#1}{#2}%
  }%
}

% Variante en largeur reduite, meme robustesse
\newcommand{\capturepretedemi}[2]{%
  \IfFileExists{#1}{%
    \begin{figure}[H]
      \centering
      \includegraphics[width=0.60\textwidth]{#1}
      \caption{#2}
    \end{figure}
  }{%
    \captureaplacer{#1}{#2}%
  }%
}

% ===========================================================================
%  ENCADRES THEMATIQUES
% ===========================================================================

% --- Encadre « decision d'architecture »
\newenvironment{decision}[1]{%
  \begin{framed}
  \noindent\textcolor{bleufonce}{\textbf{Décision --- #1}}\par\vspace{2pt}
  \small
}{%
  \end{framed}
}

% --- Encadre « constat / point d'attention »
\newenvironment{constat}[1]{%
  \begin{framed}
  \noindent\textcolor{orangeattention}{\textbf{Constat --- #1}}\par\vspace{2pt}
  \small
}{%
  \end{framed}
}

% --- Encadre « verification effectuee »
\newenvironment{verification}[1]{%
  \begin{framed}
  \noindent\textcolor{vertfonce}{\textbf{Vérification --- #1}}\par\vspace{2pt}
  \small
}{%
  \end{framed}
}

% ===========================================================================
%  CHAINE DE DEMONSTRATION EXIGEE PAR LE SUJET
%  Pratique -> Outil -> Integration -> Detection -> Resultat -> Correction
% ===========================================================================
\newcommand{\chainedemo}[6]{%
  \begin{center}
  \small
  \begin{tabular}{@{}>{\bfseries}p{3.1cm}p{10.4cm}@{}}
    \toprule
    Pratique      & #1 \\
    Outil         & #2 \\
    Intégration   & #3 \\
    Détection     & #4 \\
    Résultat      & #5 \\
    Correction    & #6 \\
    \bottomrule
  \end{tabular}
  \end{center}
}

% ===========================================================================
%  DIVERS
% ===========================================================================
\setlist[itemize]{itemsep=2pt,topsep=4pt,parsep=0pt}
\setlist[enumerate]{itemsep=2pt,topsep=4pt,parsep=0pt}

\captionsetup{font=small,labelfont=bf}

% Eviter les coupures de mots trop agressives
\tolerance=1000
\emergencystretch=2em

% Profondeur de la table des matieres
\setcounter{tocdepth}{2}
\setcounter{secnumdepth}{3}
